\chapter{电厂锅炉运行}

\section{思考题}

\begin{problem}
何谓锅炉的热平衡及其目的是什么？锅炉中存在哪些热损失？各是如何形成的？
\end{problem}

\begin{solution}
\begin{enumerate}
  \item \textbf{锅炉热平衡}：指送入锅炉的热量等于锅炉的有效利用热量与各项热损失之和。
  \item \textbf{目的}：为了评定锅炉的热经济性。
  \item \textbf{热损失及形成原因}：
  \begin{enumerate}
    \item \textbf{排烟热损失 (\(q_{2}\))}：由排烟直接带走的热量形成，是现代电厂锅炉热损失中最大的一项。
    \item \textbf{化学不完全燃烧热损失 (\(q_{3}\))}：由于烟气中含有可燃气体（如一氧化碳、氢、甲烷等）最终未能发生燃烧而造成。
    \item \textbf{机械不完全燃烧损失 (\(q_{4}\))}：由于排出的灰渣和飞灰中含有未燃尽的残炭造成。
    \item \textbf{散热损失 (\(q_{5}\))}：锅炉在运行中，汽包、联箱、汽水管道、炉墙等的温度高于外界空气温度，部分热量散失到空气中形成。
    \item \textbf{灰渣物理热损失 (\(q_{6}\))}：锅炉炉渣排出炉外时带出的高温热量形成。
  \end{enumerate}
\end{enumerate}
\end{solution}

\begin{problem}
锅炉运行调节的主要任务是什么？
\end{problem}

\begin{solution}
\begin{enumerate}
  \item 使蒸发量适应外界负荷的需要。
  \item 保证输出蒸汽的品质（蒸汽压力、温度等）。
  \item 维持正常的汽包水位。
  \item 维持高效率的燃烧和传热。
  \item 保证设备长期安全经济运行。
\end{enumerate}
\end{solution}

\begin{problem}
如何进行蒸汽压力的调节？蒸汽压力调节的一般方法是怎样的？
\end{problem}

\begin{solution}
\begin{enumerate}
  \item \textbf{调节原理}：蒸汽压力的变化实质上是锅炉蒸发量与外界负荷间的平衡被破坏，因此蒸汽压力调节实际上就是燃料量与风量的调节。无论何种扰动引起汽压变化，都应改变燃煤量和送风量，同时兼顾汽包水位及蒸汽温度的调节。
  \item \textbf{一般方法}：
  \begin{enumerate}
    \item \textbf{负荷增加，汽压下降时}：在水位正常的情况下，先与引风机配合增大送风量，再增大燃料量，强化燃烧以提高蒸发量。
    \item \textbf{负荷减少，汽压升高时}：如果此时水位较高，应先减少燃料量和送风量以减弱燃烧，再适当减少给水量或暂停给水，待汽压和水位稳定后再按正常情况调整。
  \end{enumerate}
\end{enumerate}
\end{solution}

\begin{problem}
为什么要进行蒸汽温度的调节？其调节方法一般有哪几种？目前大容量锅炉中常采用的汽温调节方法是什么？
\end{problem}

\begin{solution}
\begin{enumerate}
  \item \textbf{调节原因}：维持稳定的汽温是保证机组安全和经济运行所必需的。汽温过高会超过设备允许温度，加速钢材蠕变，严重超温甚至导致管子过热爆破。汽温过低除降低经济性外，还会增加汽轮机尾级蒸汽湿度，加剧叶片侵蚀，严重时引发水冲击。
  \item \textbf{调节方法分类}：
  \begin{enumerate}
    \item \textbf{蒸汽侧调节}：通过表面式减温器或喷水式减温器改变蒸汽温度。
    \item \textbf{烟气侧调节}：通过改变火焰中心位置、引入烟气再循环或使用烟气挡板来改变烟气侧换热。
  \end{enumerate}
  \item \textbf{大容量锅炉常采用的方法}：通常利用摆动式燃烧器改变火焰中心位置作为蒸汽温度的粗调手段，而细调则主要采用喷水减温器。
\end{enumerate}
\end{solution}

\begin{problem}
汽包水位变化过大会对锅炉运行带来哪些危害？何谓虚假水位？其影响是怎样的？
\end{problem}

\begin{solution}
\begin{enumerate}
  \item \textbf{高水位的危害}：汽水分离空间减小导致蒸汽带水，使蒸汽品质恶化。严重满水时含盐蒸汽会导致过热器结垢、超温爆管，并引发主蒸汽管道和汽轮机水冲击。
  \item \textbf{低水位的危害}：可能导致下降管带汽，使自然水循环的流动压头减小、安全性降低。若给水中断，几十秒内即可引发严重的"干锅"事故。
  \item \textbf{何谓虚假水位}：当负荷变化较大时，由于锅筒内汽、水压力瞬间不平衡而产生的水位反常现象。例如负荷突增导致汽压下降，水位会先上升后下降；负荷突降时，水位会先下降后上升。
  \item \textbf{影响}：极易干扰运行人员对实际水量的判断，如果未能正确识别虚假水位，会导致给水操作误判和误动作。
\end{enumerate}
\end{solution}

\begin{problem}
如何判断炉内燃烧工况是否正常稳定？
\end{problem}

\begin{solution}
\begin{enumerate}
  \item 当燃料量、送风量和引风量配合较好时，炉膛内应具有光亮的金黄色火焰，且火焰中无明显星点。
  \item 火焰中心应位于炉膛中部，均匀地充满整个炉膛而不触及四周的水冷壁。
  \item 烟囱排放出的烟气呈淡灰色。
\end{enumerate}
\end{solution}
